\section{Torus knots in the fundamental representation: the positivity criterion}\label{sec:torus}

\subsection{Input and the \texorpdfstring{$h/g$}{h/g} decomposition}
The input is the general formula of \S6 of \cite{AgSc12} for $T(2,2m+1)$ in the fundamental representation (normalized by the unknot); for $m = 1$ it is (6.15) there.
Under the identifications $\SO(4) = \mathrm{SU}(2)^2$ and $\SO(6) = \mathrm{SU}(4)$, (6.15) agrees with the $\U(n)$ values obtained from the construction of \cite{GN13} (checked exactly at two rational points).
At $t = q$ it agrees with the mirror image of the Kauffman polynomial of \cite{Marino09} up to a framing monomial.

% [T1]
\begin{hypothesis}[$h/g$ decomposition]\label{hyp:hg}
Write $f_{\yone}(K) = -W_{\SO}(K) = \tfrac12 h_{\yone} + g_{\yone}$ with the cover part
\[
\tfrac12 h_{\yone} = \mathcal{P}(\ba_{\mathrm{KN}}, \bq, \bt)\,\frac{\lambda - \lambda^{-1}}{\bq - \bq^{-1}},\qquad \ba_{\mathrm{KN}}^2 = a\,(t/q)^{1/2},
\]
where $\mathcal P$ is the superpolynomial of \cite{DGR05}, i.e., the A-type refined Chern--Simons invariant in the fundamental representation.
\end{hypothesis}

% [T2]
\begin{result}\label{cr:monomial}
For $m = 1, \dots, 5$ we searched for monomials $\pm q^{k_1/2} t^{k_2/2} (t^N)^{k_3}$ for which $g_{\yone}$ becomes a Laurent polynomial. The only solutions have $k_3 = -2(2m+1)$ and $k_1 = -(3m+2)$; $k_2$ is fixed uniquely to $k_2 = 7m+4$ by requiring the $t = q$ limit to have the form of a Kauffman polynomial. The monomial is therefore $T_{\yone}^{-2(2m+1)}(q/t)^{\delta/2}$ with $T_{\yone} = q^{1/2}t^{N-1}$ and $\delta = m$.
In the uniform normalization the extra $(q/t)^{\delta/2}$ moves to the cover, and the $\SO$ monomial is exactly $T_{\yone}^{-2(2m+1)}$ (full agreement with the data at fixed $N = 2, 3$ for $m \le 3$).
\end{result}
The same factor $\sqrt{q/t}$ appears for the $\mathrm{SU}(N)$ trefoil in \cite{AgSh11}.

\subsection{Results}
Write $g^{\mathrm{true}}_{\yone} = (t/q)^{\delta/2} g_{\yone}$ for the crosscap part in the uniform normalization.
% [T3][T4]
\begin{result}\label{cr:torus}
For $m = 1, \dots, 5$, $g^{\mathrm{true}}_{\yone}$ is an integer Laurent polynomial (with 11, 27, 49, 77, 111 terms) and:
\begin{enumerate}
\item[(N)] $g^{\mathrm{true}}_{\yone} = 1$ at the $\SO(2)$ point $\ba^2 = \bq^2$ ($\lambda^2 = t$).
\item[(P)] It is invariant under $\bt \to -\bt$.
\item[($c{=}2$)] The part with odd powers of $\ba$ is
\begin{equation}\label{eq:Pim}
-(\ba^2 - \bq^2)(\bt^2 - 1)\,\ba^{-(4m+1)}\bq^{-(2m-1)}\bt^{-4m}\,\Pi_m,\qquad
\Pi_m = \sum_{j=0}^{m-1}\ba^{2j}\bq^{2(m-1-j)}\bt^{2j}\sum_{i=0}^{j}(\bq^4\bt^2)^i ,
\end{equation}
where $\Pi_m$ has $m(m+1)/2$ terms, all with coefficient $+1$.
\item[($c{=}1$)] The part with integer powers of $\ql$ (even powers of $\ql^{1/2}$) is $\sigma$-symmetric and is a polynomial in $G$ with coefficients in $\mathbb{Z}[\qr^{\pm1}]$. Each $\lambda^\beta$ block is one of
  \begin{itemize}
  \item $\lambda^{-(4m+2)}$: a positive monomial;
  \item $\lambda^{-4m}, \dots, \lambda^{-(2m+4)}$: $(\bt^2 - 1)\times$ positive;
  \item $\lambda^{-(2m+2)}$: all coefficients negative;
  \item $\lambda^{-2m}$: all coefficients positive.
  \end{itemize}
\end{enumerate}
\end{result}
\begin{conjecture}\label{conj:Pim}
The closed form \eqref{eq:Pim} of the $c = 2$ part holds for all $m$. (It was guessed from $m \le 3$ and confirmed for $m = 4, 5$.)
\end{conjecture}
At $t = q$ the odd powers of $z$ disappear, consistent with the vanishing of the unrefined $c = 2$ invariants for torus knots in the fundamental representation. The values for $m = 1, 2, 3$ are listed in Appendix \ref{app:torus}.

\subsection{The criterion}\label{sec:criterion}
The unrefined crosscap BPS numbers are integers whose sign depends on $\beta$ (e.g., $\hat g_{\yone,\yone} = z(\lambda - \lambda^{-1})$ for the Hopf link in \cite{Marino09}).
After refinement one expects positivity, which we formulate as follows: in each $\lambda^\beta$ block the coefficients have a common sign, equal to the unrefined sign $\varepsilon_\beta$.

\begin{definition}\label{def:c1}
Expand the $c = 1$ part of the crosscap invariant in the fundamental representation as
\[
g^{c=1} = \sum_\beta \lambda^\beta \sum_{g, J} N^{c=1}_{g,\beta,J}\, G^g \qr^{2J},\qquad 2J \in \mathbb{Z}.
\]
At $t = q$, $\sum_J N^{c=1}_{g,\beta,J}$ reduces to the $N^{c=1}_{g,\beta}$ of \cite{Marino09}.
\end{definition}
For a knot whose unrefined $c = 2$ invariants vanish, such as a torus knot in the fundamental representation, the whole $c = 2$ part is ``refined'' in the sense of \S\ref{sec:universal}; in general only the refined part $X_1^{c=2} = g^{c=2} - g^{c=2}|_{t=q}(\lambda \to \tlam)$ can be expected to carry the universal factor.

% [T6]
\begin{conjecture}[criterion]\label{conj:criterion}
For two-strand torus knots $T(2,2m+1)$ in the fundamental representation:
\begin{enumerate}
\item[($c{=}1$)] after the $G$-expansion of Definition \ref{def:c1}, each $\lambda^\beta$ block has a common sign, or is $(\bt^2 - 1)\times$ a block with a common sign, the sign being the unrefined sign $\varepsilon_\beta$;
\item[($c{=}2$)] the refined part of the $c = 2$ part is $(\ba^2 - \bq^2)(\bt^2 - 1)\times$ a polynomial whose coefficients have a common sign in each $\ba$ block; for $T(2,2m+1)$ it is the whole $c = 2$ part, $-(\ba^2 - \bq^2)(\bt^2 - 1)\times$ a positive polynomial.
\end{enumerate}
Items (N) and (P) of Result \ref{cr:torus} are checks of the normalization; the content on positivity is in ($c{=}1$) and ($c{=}2$).
\end{conjecture}
The evidence is $T(2,2m+1)$, $m \le 5$ (Result \ref{cr:torus}). We do not state the criterion for general knots: for $T(3,4)$ the $c = 1$ part is not $\sigma$-symmetric, so ($c{=}1$) fails while ($c{=}2$) holds; for the algebraic cables of \S\ref{sec:cable} the $c = 1$ part is not $\sigma$-symmetric and the refined part of the $c = 2$ part has mixed signs within $\ba$ blocks (Result \ref{cr:cable-largeN}). Whether the Kauffman superpolynomials of \cite{GW05} for two-strand torus knots, which are positive, are related to our $c = 1$ part is not examined here \open.
The factor $(\ba^2 - \bq^2)$ vanishes at the $\SO(2)$ point, where knot invariants of an abelian group become trivial; $(\bt^2 - 1)$ follows from vanishing at $t = q$ together with the parity (P).
